of a projective morphism of complex analytic spaces, and will undoubtedly also apply in arbitrary characteristic once the problem of the ``exterior power'' operations in the derived category has been solved (cf.\ Exp.\ XIV, n\textsuperscript{os}~1, 2, 3). The absence of a study of this question undoubtedly constitutes the most serious gap in this Seminar, even from the point of view we have adopted.

One will notice the absence, from the table of contents of the present Seminar, of two of the participants mentioned on the cover page. One of them, J.P. Jouanolou, had taken an active part in the technical development of the first part of the seminar, but had been prevented from taking part in the oral exposés; to him we owe, in particular, the linear independence assertion contained in the important statement VI 1.1. giving the structure of the ring @@M0@@ of classes of locally free Modules on a projective bundle (which had previously been proved only under the assumption that the base scheme admits an ample invertible Module). The other, J.P. Serre, had given two oral exposés, logically independent of the rest of the Seminar, which he therefore preferred to publish as a separate article (cf.\ J.P. Serre, Groupes de Grothendieck des schémas en groupes réductifs déployés, to appear in Publications Mathématiques, n\textsuperscript{o}~34).

As in the other parts of the Algebraic Geometry Seminar of Bois Marie, the sigla SGA 1 to SGA 6 refer to the different parts of that Seminar, the Roman numeral following the siglum indicating the number of the exposé, and the Arabic numerals that follow it corresponding to the internal numbering of the exposé in question; for internal references within the present Seminar, the siglum SGA 6 is omitted, using references of the form I 4.9 (meaning: Exposé I, statement 4.9). The siglum EGA refers to the Éléments de Géométrie Algébrique of J. Dieudonné and A. Grothendieck.

